% ==============================================================================
% 05_discussion.tex – Chapter 5
% ==============================================================================
\chapter{Discussion}
\label{ch:diskussion}

\section{Answers to the Research Questions}
\label{sec:antworten}

\paragraph{RQ1: How much does a short personalization help?}
The generic base model is not better than the constant baseline on the own
recordings (\SI{30.57}{\percent} versus \SI{30.32}{\percent}). With a temporal
split, which matches real use, personalization improves the result. The error
falls from \SI{31.84}{\percent} to \SI{27.93}{\percent}. This is an improvement
of about \SI{12}{\percent} (Wilcoxon $p=0.002$). In 8 of 18 sessions, however,
there is no improvement (\cref{sec:design_performance}). In the good sessions 21
and 22, the improvement of \SI{46}{\percent} is still clear. For a second
recording, too, the gain stays almost the same at \SI{47}{\percent}
(\cref{sec:loading}).

With a random split, the results are even better. The error falls to
\SI{19.79}{\percent}, i.e.\ by about one third. The best session even reaches an
error of only \SI{10.5}{\percent} (\cref{sec:sessions_results}). This split,
however, shows a result that is too positive and does not match real use.

For person~B, personalization with the whole camera image brought no
improvement in either session. This also holds for the random split. With the
eye crop of the web demo, in contrast, the error fell from \SI{26.9}{\percent}
to \SI{9.7}{\percent} (\cref{sec:person_b}).

\emph{Answer:} a short personalization can reduce the error clearly. It does
not, however, work equally well for every recording. Good recording conditions
and a suitable image crop of the eyes are important for it. A random split also
shows better results than one can expect in real use.

\paragraph{RQ2: Which training setting is suitable?}
On GazeCapture, most training settings show no clear difference. The results of
all models are close to the constant baseline (\cref{sec:training}). The tests
show, however, that sigmoid as the activation gives the best results. In
addition, the L1 loss is better than the loss with position weighting. It is
enough to train only the regression head. This gives similar results to
training further layers, but is more stable and much faster.

\emph{Answer:} the following setting is suitable for training:
\begin{itemize}[nosep]
  \item frozen backbone (not trained),
  \item regression head with sigmoid,
  \item L1 loss,
  \item learning rate $10^{-4}$,
  \item batch size 32.
\end{itemize}
This setting was chosen mainly because of its stability and the low training
effort. The learning rate $10^{-4}$ was also chosen because it reached the
lowest error in the first test.

\paragraph{RQ3: What is the role of data split, calibration and smoothing?}
The data split has a large influence on the results. A random split gives
results that are too good, because almost identical frames of one recording
appear both in training and in test. The smooth movement of the target over
the screen (Gen~1) covers the screen best and reaches the lowest error; the
difference, however, is not significant (\cref{sec:design_performance}). Fast
jumps produce frames with wrong labels. Removing these WAIT frames improves the
result slightly, by \num{0.84} percentage points. This improvement, however, is
not significant. More data in the corners does not improve the results in these
areas. A moving average over the predictions makes the display much calmer.
With a window of $k=3$, the RMS-S2S falls by \SI{40}{\percent}. The accuracy,
however, does not improve. With longer windows, the result even gets worse,
because the prediction is delayed more (\cref{sec:smoothing}).

\emph{Answer:} for calibration, a short and smooth movement over the whole
screen is best. Fast jumps should be detected and logged, so that the frames
during these movements can be removed later. For evaluation, a temporal split
should always be used, because it matches real use better. A short moving
average (about \SI{0.5}{\second}) can make the display calmer. It does not,
however, improve the accuracy and does not replace a better calibration.

\paragraph{RQ4: What is the influence of the recording conditions?}
The largest influence comes from the background. A busy background raises the
error by 6 to 15 percentage points. Glasses also lead to a higher error in most
cases. In dark rooms, personalization helps only a little. The length of the
recording and the device used (webcam), in contrast, show no clear influence
(\cref{sec:conditions}).

\emph{Answer:} a calm background and good lighting are particularly important
for good results. They have a larger influence than the type of calibration.
These results, however, are based only on single sessions and must therefore be
read with care.

\section{Relation to the State of the Art}
\label{sec:einordnung}
Some results fit the literature well:
\begin{itemize}
  \item \textbf{Few parameters are enough.} As in offset calibration
        \cite{chen2020offset}, SPAZE \cite{linden2019spaze} and PIGaze
        \cite{wang2026pigaze}, a small adapted head on a frozen backbone is
        enough here. More trainable layers only make training less stable. This
        fits Kumar et al.~\cite{kumar2022finetuning}.
  \item \textbf{Personalization is the most important factor.} As in Valliappan
        et al.~\cite{valliappan2020smartphone}, almost the whole gain comes from
        a few seconds of personal data.
  \item \textbf{The evaluation protocol decides.} Similar to the difference
        between \code{norm\_random} and \code{norm\_subject}, random splits
        overestimate the performance (\cref{subsec:norm_subject}).
  \item \textbf{Smoothing improves precision, not accuracy.} This matches the
        distinction by Niehorster et al.~\cite{niehorster2026quality}: a moving
        average clearly reduces the scatter from frame to frame, but the
        systematic error stays almost the same (\cref{sec:smoothing}).
  \item \textbf{Wrong labels have a small effect.} DMRF \cite{jiang2026dmrf} also
        lowers the error only a little ($3.92^\circ \rightarrow 3.85^\circ$).
        Here, removing the WAIT frames gives \num{0.84} percentage points;
        this difference is not significant ($p=0.43$).
\end{itemize}
There is also an important difference: the generic model of this thesis is very
weak. It stays at the screen centre, while methods in the literature reach a few
degrees of accuracy even without personalization. The main reason is probably the
input: most methods use face or eye crops, while this thesis uses the whole
camera image. In addition, the base model uses only \num{1000} frames of 179
persons. The absolute numbers are in any case not directly comparable with the
literature because of the different metrics (\cref{sec:vergleich}).

\section{Limitations}
\label{sec:grenzen}
\begin{itemize}
  \item \textbf{Few persons:} the 18 evaluated sessions come from one person
        (the author). From a second person there are only two sessions and one
        run of the demo. That the eye crop explains the difference is a hint
        from this single run and was not examined systematically. The results
        therefore do not simply transfer to other persons.
  \item \textbf{Few sessions and mixed conditions:} 18 sessions are not enough
        for reliable statements about single factors. Design, light, background
        and glasses often change at the same time.
  \item \textbf{Single runs:} there is usually only one run per session (seed
        42). The spread between seeds can be up to 3 percentage points.
  \item \textbf{One backbone and whole images:} only ResNet-18 with whole camera
        images was studied. DINOv2, CLIP, MAE, UniGaze, adapters and LoRA were
        not tested, and neither were face or eye crops in the experiments of
        Chapter~4.
  \item \textbf{Weak base model:} the base model hardly learns anything on
        GazeCapture and mainly serves as a starting point for personalization
        in this thesis. Whether a randomly initialized head would work equally
        well was not tested. A better base model could change the
        personalization results.
  \item \textbf{Many tests:} many statistical tests were computed without a
        correction for multiple comparisons (for example Bonferroni). Close
        results such as $p=0.032$ (L1 versus the sensitive loss) should
        therefore be read with care.
  \item \textbf{Smoothing only afterwards:} the moving average was evaluated
        on stored predictions of the sequential split (mostly one seed) and of
        the two videos. For fixations, there is only a single run of the demo.
  \item \textbf{Metric:} the diagonal error is not directly comparable with
        angular errors or centimetres from the literature.
\end{itemize}
